\documentclass[conference]{IEEEtran}
\IEEEoverridecommandlockouts
\usepackage{cite}
\usepackage{amsmath,amssymb,amsfonts}
\usepackage{algorithmic}
\usepackage{graphicx}
\usepackage{textcomp}
\usepackage{xcolor}
\usepackage{booktabs}
\usepackage{url}
\def\BibTeX{{\rm B\kern-.05em{\sc i\kern-.025em b}\kern-.08em
    T\kern-.1667em\lower.7ex\hbox{E}\kern-.125emX}}

\begin{document}

\title{Agentic Autoresearch for Discrete Text
Diffusion Transformers on Consumer Hardware\\
A Controlled Comparison of Uniform (SEDD) and Masked (MDLM) Objectives}

\author{\IEEEauthorblockN{Goutham Krishnan}
\IEEEauthorblockA{goutham.krishnan22@gmail.com}}

\maketitle

\begin{abstract}
Discrete diffusion has emerged as a non-autoregressive alternative to
left-to-right language modeling, but most reported results depend on
large compute budgets that place careful architectural exploration out
of reach for practitioners on consumer hardware. We present TON-v1, a
77.2M-parameter bidirectional diffusion transformer trained to generate
short children's stories on a single 8\,GB RTX 4060 laptop GPU. Our
central methodological contribution is \emph{agentic autoresearch}: an
autonomous large-language-model agent iteratively edits the training
script, trains under a fixed wall-clock budget, reads back a target
metric, and keeps an edit only if the metric improves. Using this loop
we conduct two controlled studies on an identical transformer backbone.
The first tunes a Score Entropy Discrete Diffusion (SEDD, uniform)
model over 52 experiments, reducing validation loss by 39\%, with a
single change, replacing learned absolute position embeddings with
rotary embeddings (RoPE), accounting for a 28\% reduction on its own.
The second replaces the diffusion objective with a Masked Diffusion
Language Model (MDLM, absorbing) formulation and, guided by a
loss-agnostic generative-perplexity metric, reduces generative
perplexity from 350.5 to 42.6 over 12 experiments (an
$8.2\times$ improvement). Full training of the best MDLM configuration
for 100{,}000 steps reaches a generative perplexity of 24.75 with a
distinct-2 diversity of 0.68, producing fluent, non-degenerate text.
We report full ablation logs, generation-time latency and memory
benchmarks, and a candid account of which interventions helped and
which did not, as a reproducible baseline for small-scale discrete
diffusion research.
\end{abstract}

\begin{IEEEkeywords}
discrete diffusion, masked diffusion, text generation, transformers,
non-autoregressive generation, automated machine learning, rotary
position embeddings, consumer GPU
\end{IEEEkeywords}

\section{Introduction}
Autoregressive (AR) language models generate text one token at a time,
conditioning each token only on its predecessors. Discrete diffusion
models offer an alternative: they generate an entire sequence in
parallel by iteratively denoising a corrupted sequence, allowing every
position to attend bidirectionally to every other at every step. Recent
formulations, Score Entropy Discrete Diffusion (SEDD)~\cite{sedd} and
Masked Diffusion Language Models (MDLM)~\cite{mdlm}, have narrowed the
perplexity gap to AR models, but the design decisions behind these
results are typically established at a scale unavailable to most
practitioners.

This paper asks a practical question: \emph{how far can a small discrete
diffusion transformer be pushed on a single consumer GPU, and how much
of the tuning can be delegated to an autonomous agent?} We make three
contributions:

\begin{enumerate}
\item \textbf{An agentic autoresearch methodology} in which an LLM agent
autonomously proposes, implements, runs, and accepts-or-rejects
training-script edits under a fixed wall-clock budget, using an
objective metric as the sole acceptance criterion.
\item \textbf{A controlled comparison} of the uniform (SEDD) and
masked/absorbing (MDLM) diffusion objectives on an \emph{identical}
77.2M-parameter backbone, isolating the effect of the corruption process
and loss from confounding architectural differences.
\item \textbf{A fully reproducible baseline}, including complete
per-experiment ablation logs, generation-quality metrics, and
inference-time latency/memory benchmarks, all obtained on an 8\,GB RTX
4060 laptop GPU.
\end{enumerate}

\section{Background and Related Work}

\subsection{Discrete Diffusion}
Continuous diffusion models~\cite{sohl,ddpm} corrupt data with Gaussian
noise and learn to reverse the process. For discrete data such as text,
the corruption is instead a Markov chain over token states~\cite{d3pm}.
Two corruption kernels are relevant here. In the \emph{uniform} kernel,
a token is replaced by another token drawn uniformly from the
vocabulary; SEDD~\cite{sedd} learns the concrete score (the ratio of
marginal probabilities between states) and optimizes a denoising score
entropy objective. In the \emph{absorbing} kernel, a token is replaced
by a dedicated \texttt{[MASK]} state that is never reversed except by
prediction; MDLM~\cite{mdlm} shows that this reduces to a simple,
time-weighted masked cross-entropy objective, essentially a
diffusion-theoretic generalization of masked language modeling
\cite{bert}, and reports improved perplexity over SEDD.

\subsection{Architectural Ingredients}
Our backbone is a standard transformer encoder~\cite{transformer} with
bidirectional attention. We adopt rotary position embeddings
(RoPE)~\cite{rope} for relative positional encoding, adaptive layer
normalization (adaLN)~\cite{dit} to condition every block on the
diffusion noise level, and query-key normalization for attention
stability. The vocabulary and tokenizer are inherited from
GPT-2~\cite{gpt2}.

\subsection{Automated and Agentic Model Search}
Classical neural architecture search automates discrete design choices
but is compute-intensive. Recent work explores using LLM agents as
autonomous researchers that write and evaluate their own code. Our loop
follows the spirit of Karpathy's ``autoresearch''
formulation~\cite{autoresearch}: a tight edit-train-measure-decide
cycle on a single script under a fixed budget. To our knowledge this is
among the first documented applications of such a loop to discrete
diffusion language models on consumer hardware.

\section{Model Architecture}
Both studies share one backbone (Table~\ref{tab:arch}). A sequence of
$T{=}128$ token embeddings of width $d{=}512$ passes through six
identical pre-norm transformer blocks. Each block applies
(i)~multi-head self-attention with 16 heads, RoPE, and RMSNorm applied
to per-head queries and keys, followed by (ii)~a position-wise
feed-forward network of hidden width $4d$. Attention is computed with a
fused scaled-dot-product kernel and is fully bidirectional. The
diffusion time $t\in[0,1]$ is mapped through a sinusoidal embedding and
an MLP to a conditioning vector that modulates each block via adaLN
(a per-block scale and shift on the normalized activations).

The only architectural difference between the two models is the token
embedding table. The MDLM variant adds one row for the absorbing
\texttt{[MASK]} state, giving $|\mathcal{V}|{+}1$ input rows against a
$|\mathcal{V}|$-way output head; the SEDD variant uses
$|\mathcal{V}|$ rows throughout. This accounts for a difference of
exactly $d{=}512$ parameters (77{,}236{,}689 for SEDD versus
77{,}237{,}201 for MDLM), i.e. both models are $\approx$77.2M
parameters.

\begin{table}[t]
\caption{Shared Backbone and Training Configuration}
\label{tab:arch}
\centering
\begin{tabular}{@{}ll@{}}
\toprule
\textbf{Component} & \textbf{Value} \\
\midrule
Parameters & $\approx$77.2\,M \\
Layers / width / heads & 6 / 512 / 16 \\
Context length & 128 tokens \\
Attention & bidirectional, RoPE (base 1000), QK-norm \\
Time conditioning & sinusoidal $\to$ MLP $\to$ adaLN per block \\
Tokenizer / vocabulary & GPT-2 BPE / 50{,}257 ($+1$ for MDLM) \\
Precision & bf16 autocast, TF32 matmuls \\
Optimizer & fused Adam, grad-clip 1.0 \\
Batch size & 16 \\
LR schedule & 100-step warmup, cosine decay \\
Peak LR (SEDD / MDLM) & $1\times10^{-3}$ / $3\times10^{-3}$ \\
\bottomrule
\end{tabular}
\end{table}

\section{Training Objectives}
Let $x_0$ be a clean token sequence and $t\sim\mathcal{U}(0,1)$ a
sampled noise level (drawn with stratified sampling across the batch to
reduce gradient variance).

\subsection{SEDD (Uniform)}
The forward process replaces tokens with uniform random tokens under a
geometric noise schedule. The model outputs per-position log-ratios
$s_\theta$, trained with a denoising weighted score entropy (DWDSE)
objective. We use a closed-form evaluation of the loss that avoids
materializing the full rate and transition tensors, which reduced the
forward pass from 72\,ms to 35\,ms per step in our implementation.

\subsection{MDLM (Absorbing)}
The forward process independently replaces each token with
\texttt{[MASK]} with probability $t$ (a linear schedule). The model
predicts the original token at each masked position, and the negative
ELBO reduces to a time-weighted cross-entropy over masked positions
only:
\begin{equation}
\mathcal{L}_{\text{MDLM}} =
\mathbb{E}_{t, x_0, x_t}\!\left[
\frac{1}{t}\sum_{i:\,x_t^{(i)}=\texttt{[MASK]}}
-\log p_\theta\!\big(x_0^{(i)} \mid x_t\big)
\right].
\end{equation}
The $1/t$ weight upweights lightly-corrupted examples, whose correct
reconstruction is a stronger learning signal. Generation runs the
reverse process from an all-\texttt{[MASK]} sequence, progressively
unmasking positions over $N$ steps; at each step we sample predicted
tokens with top-$k$ filtering.

\section{Agentic Autoresearch Methodology}
The tuning of both models was performed by an autonomous LLM agent, not
by hand. The agent executes the following loop indefinitely until
stopped:

\begin{enumerate}
\item Inspect the current git state and the training script.
\item Propose and directly implement one experimental change (any part
of the architecture, optimizer, schedule, loss, or sampler is in scope).
\item Commit, then launch training under a fixed wall-clock budget
(5 minutes for SEDD, 10 minutes for MDLM), redirecting all output to a
log.
\item Parse the final target metric from the log.
\item If the metric improved, advance the branch (keep the commit);
otherwise revert to the prior state. Log the outcome, with a one-line
rationale, to a tab-separated results file.
\end{enumerate}

Because the budget is wall-clock rather than a fixed step count, the
loop implicitly rewards throughput: any change that increases
steps-per-second lets the step-bound model train longer within the same
budget, which is itself a source of metric improvement. This design
made low-level efficiency work (mixed precision, fused kernels,
closed-form losses) a first-class part of the search rather than an
afterthought.

\subsection{Choice of Target Metric}
The two objectives are not numerically comparable: DWDSE and weighted
cross-entropy live on different scales, and the raw SEDD loss values we
report are unnormalized, implementation-specific quantities useful only
for \emph{relative} comparison within a single study. For the SEDD study
we therefore optimized held-out DWDSE directly. For the MDLM study,
where the goal was to compare against and surpass the SEDD baseline, we
introduced a loss-agnostic quality metric: \textbf{generative
perplexity} of the model's own samples under a frozen
GPT-2~\cite{gpt2} scorer (valid because our samples are already GPT-2
tokens), guarded by a \textbf{distinct-2} bigram-diversity measure to
detect degenerate low-perplexity output (e.g., repetition). All MDLM
acceptance decisions used generative perplexity, with distinct-2 as a
veto.

\section{Experimental Setup}
All experiments ran on a single NVIDIA RTX 4060 Laptop GPU (8\,GB) under
Windows. Data is the \texttt{karpathy/tinystories-gpt4-clean}
corpus~\cite{tinystories}, a set of simple short stories; the full
corpus of $\approx$542M GPT-2 tokens was tokenized once and cached to
disk, with a 90/10 train/validation split. The autoresearch loops used
a token subset for speed; the final full-training runs used the entire
corpus. Held-out evaluation uses a fixed random seed so that reported
losses depend only on model weights and are comparable across commits
within a study.

\section{Results}

\subsection{SEDD Autoresearch}
Over 52 experiments (22 accepted), the agent reduced held-out DWDSE from
a baseline of 325{,}306 to 198{,}472, a 39\% reduction
(Table~\ref{tab:sedd}). Early gains were dominated by throughput
interventions: TF32 matmuls, fused QKV with scaled-dot-product
attention, a training-only bf16 pass, and the closed-form DWDSE
loss, each of which increased steps-per-second and thus effective
training within the fixed budget. A sequence of small architectural
wins followed: QK-norm, sinusoidal time conditioning, adaLN, 16 attention
heads, and a reduction from 8 to 6 layers (the model is step-bound, not
capacity-bound, so depth traded unfavorably against steps). The single
largest intervention was replacing the learned absolute position
embedding with RoPE, which reduced the loss by 28\% in one step
(279{,}226 $\to$ 201{,}269); tuning the RoPE base to 1000 for the short
context yielded the best result of 198{,}472. The agent also correctly
rejected numerous plausible-but-unhelpful ideas, including weight tying
(which diverged, as the output head feeds an exponential and cannot
share the input embedding), SwiGLU, wider embeddings, importance
sampling of $t$, and both larger and smaller batch sizes.

\begin{table}[t]
\caption{Selected SEDD Autoresearch Milestones (Held-out DWDSE)}
\label{tab:sedd}
\centering
\begin{tabular}{@{}lrl@{}}
\toprule
\textbf{Intervention} & \textbf{Val loss} & \textbf{Note} \\
\midrule
Baseline & 325{,}306 & tuned starting point \\
TF32 matmul & 290{,}336 & throughput \\
Fused Adam & 294{,}316 & throughput \\
QK-norm & 286{,}682 & stability \\
Closed-form DWDSE & 285{,}160 & +27\% steps/s \\
16 heads & 284{,}253 & finer attention \\
adaLN (scale+shift) & 280{,}337 & time conditioning \\
\textbf{RoPE} & \textbf{201{,}269} & \textbf{single largest win} \\
RoPE base 1000 & \textbf{198{,}472} & best \\
\bottomrule
\end{tabular}
\end{table}

\subsection{MDLM Autoresearch}
Starting from the tuned SEDD backbone, the agent replaced the objective
with MDLM and, guided by generative perplexity, ran 12 experiments
(7 accepted). Switching the corruption process and loss alone, with no
architectural change, reduced generative perplexity from 350.5 to
221.7 in a single step while simultaneously improving diversity
(distinct-2 $0.70 \to 0.85$), confirming the reported advantage of the
absorbing objective. Raising the peak learning rate to $3\times10^{-3}$,
increasing the number of sampling steps to 256 (the single largest
sampling win, $186\to115$), and top-$k{=}20$ sampling drove generative
perplexity to 42.6 (Table~\ref{tab:mdlm}). The distinct-2 guard proved
its worth: top-$k{=}10$ matched the perplexity of top-$k{=}20$ but
reduced diversity to 0.60 with visible repetition, and was correctly
rejected.

\begin{table}[t]
\caption{MDLM Autoresearch Progression (Generative Perplexity)}
\label{tab:mdlm}
\centering
\begin{tabular}{@{}lrrl@{}}
\toprule
\textbf{Intervention} & \textbf{Gen. PPL} & \textbf{Distinct-2} & \textbf{Status} \\
\midrule
SEDD baseline & 350.5 & 0.70 & keep \\
$\to$ MDLM objective & 221.7 & 0.85 & keep \\
LR $3\times10^{-3}$ & 186.2 & 0.84 & keep \\
256 sampling steps & 115.1 & 0.76 & keep \\
top-$k{=}50$ & 74.8 & 0.75 & keep \\
\textbf{top-$k{=}20$} & \textbf{42.6} & \textbf{0.66} & \textbf{keep} \\
top-$k{=}10$ & 41.8 & 0.60 & discard \\
\bottomrule
\end{tabular}
\end{table}

\subsection{Full Training}
The best configuration of each model was then trained to convergence on
the full corpus (not budget-limited). The SEDD model trained for
75{,}000 steps, reaching a best held-out DWDSE of 124{,}395. The MDLM
model trained for 100{,}000 steps, reaching a final generative
perplexity of \textbf{24.75} with distinct-2 of 0.68. For reference,
real TinyStories text scores in the approximate 30 to 50 range under the
same GPT-2 scorer; the trained MDLM model therefore produces samples
whose fluency, as measured by this proxy, is comparable to the training
distribution while retaining healthy lexical diversity. A representative
uncurated sample:

\begin{quote}\small\itshape
``One day, the smart cat saw a big dog. The dog wanted [the] tree and
the cat. The cat was in the tree\ldots\ The cat and the bird became
friends. They played and ate [by] the tree together. The cat was happy
again.''
\end{quote}

\subsection{Inference Benchmarks}
Table~\ref{tab:bench} reports single-sequence generation latency and
peak memory on the RTX 4060, averaged over 20 runs after warmup. Because
generation is non-autoregressive, there is no key-value cache and no
prompt-prefill asymmetry: every denoising step is a full-sequence
bidirectional forward pass of near-identical cost. Total latency is thus
dominated by the number of denoising steps. The MDLM model, tuned to
256 steps, is correspondingly slower per sample than the 128-step SEDD
model, but both generate a full 128-token story in a few seconds within
$\approx$1.4\,GB of memory, well inside the 8\,GB budget. We note that
per-sample latency amortizes strongly under batching, since a single
sequence underutilizes the GPU.

\begin{table}[t]
\caption{Generation Benchmarks on RTX 4060 (1 Sequence, 128 Tokens)}
\label{tab:bench}
\centering
\begin{tabular}{@{}lrr@{}}
\toprule
\textbf{Metric} & \textbf{SEDD} & \textbf{MDLM} \\
\midrule
Denoising steps & 128 & 256 \\
Per-step latency (ms) & 18.05 & 16.76 \\
Total latency (ms) & 2311 & 4293 \\
Peak memory (MB) & 1356 & 1330 \\
Cold-start (ms) & 3082 & 5303 \\
\bottomrule
\end{tabular}
\end{table}

\section{Discussion}

\subsection{What Worked}
Three themes dominated the useful interventions. First, in a step-bound
regime, \emph{throughput is quality}: mixed precision, fused kernels,
and a closed-form loss all improved the target metric purely by enabling
more optimization steps within the fixed budget. Second,
\emph{relative positional information matters disproportionately}: RoPE
was the single largest win in the SEDD study. Third, for the absorbing
objective, \emph{sampling-time decisions rival training-time ones}: the
number of denoising steps and the top-$k$ threshold together moved
generative perplexity more than any single training change after the
objective switch itself.

\subsection{What Did Not}
The agent's rejected experiments are as informative as its accepted
ones. Increased depth, increased width, and larger batch sizes all lost
to the step-bound baseline. Weight tying diverged for a structural
reason specific to the score parameterization. Importance sampling of
the noise level, SwiGLU, and RMSNorm produced no benefit at this scale.
The consistency of the ``6 layers, 512 width, 16 heads'' optimum across
many re-tests suggests these are genuine local optima for this
data-and-budget regime rather than noise.

\subsection{Limitations}
This is a small-scale study on a narrow, simple corpus; the absolute
quality is far from large diffusion or AR language models, and we make
no scaling claims. Generative perplexity under a fixed scorer is a proxy
that can be gamed by low-entropy text, which is precisely why the
distinct-2 guard was necessary, and why we caution against reading the
sub-30 figure as surpassing real text in any absolute sense. The
wall-clock budget couples results to the specific hardware. Finally, the
autoresearch loop performs greedy, single-metric hill-climbing and does
not explore interactions between rejected changes; a more sophisticated
search could recover value from combinations it discarded individually.

\section{Conclusion}
We demonstrated that an autonomous LLM agent can meaningfully tune a
discrete diffusion transformer on a single consumer GPU, and used that
loop to run a controlled comparison of the uniform (SEDD) and masked
(MDLM) diffusion objectives on an identical 77.2M-parameter backbone.
The masked/absorbing objective, combined with agent-discovered
sampling-time tuning, reduced generative perplexity by more than
$8\times$ over the tuned uniform baseline and yielded fluent,
non-degenerate short stories at a final generative perplexity of 24.75.
We release the complete ablation logs, metrics, and benchmarks as a
reproducible baseline for small-scale discrete diffusion research, and
suggest agentic autoresearch as a practical tool for architecture and
objective exploration under tight compute budgets.

\section*{Reproducibility}
All code, per-experiment logs (\texttt{results.tsv},
\texttt{results\_v2.tsv}), generated samples, loss curves, and the
benchmark harnesses used to produce every number in this paper are
organized into two self-contained experiment directories,
\texttt{diffusion\_sedd/} and \texttt{diffusion\_mdlm/}.

\begin{thebibliography}{00}
\bibitem{sohl} J. Sohl-Dickstein, E. Weiss, N. Maheswaranathan, and
S. Ganguli, ``Deep unsupervised learning using nonequilibrium
thermodynamics,'' in \emph{Proc. ICML}, 2015.
\bibitem{ddpm} J. Ho, A. Jain, and P. Abbeel, ``Denoising diffusion
probabilistic models,'' in \emph{Proc. NeurIPS}, 2020.
\bibitem{d3pm} J. Austin, D. Johnson, J. Ho, D. Tarlow, and
R. van den Berg, ``Structured denoising diffusion models in discrete
state-spaces,'' in \emph{Proc. NeurIPS}, 2021.
\bibitem{sedd} A. Lou, C. Meng, and S. Ermon, ``Discrete diffusion
modeling by estimating the ratios of the data distribution,'' in
\emph{Proc. ICML}, 2024. arXiv:2310.16834.
\bibitem{mdlm} S. Sahoo \emph{et al.}, ``Simple and effective masked
diffusion language models,'' in \emph{Proc. NeurIPS}, 2024.
arXiv:2406.07524.
\bibitem{transformer} A. Vaswani \emph{et al.}, ``Attention is all you
need,'' in \emph{Proc. NeurIPS}, 2017.
\bibitem{bert} J. Devlin, M. Chang, K. Lee, and K. Toutanova, ``BERT:
Pre-training of deep bidirectional transformers for language
understanding,'' in \emph{Proc. NAACL}, 2019.
\bibitem{rope} J. Su, Y. Lu, S. Pan, B. Wen, and Y. Liu, ``RoFormer:
Enhanced transformer with rotary position embedding,'' arXiv:2104.09864,
2021.
\bibitem{dit} W. Peebles and S. Xie, ``Scalable diffusion models with
transformers,'' in \emph{Proc. ICCV}, 2023.
\bibitem{gpt2} A. Radford \emph{et al.}, ``Language models are
unsupervised multitask learners,'' OpenAI Tech. Rep., 2019.
\bibitem{tinystories} R. Eldan and Y. Li, ``TinyStories: How small can
language models be and still speak coherent English?'' arXiv:2305.07759,
2023.
\bibitem{autoresearch} A. Karpathy, ``autoresearch,'' GitHub repository.
[Online]. Available: \url{https://github.com/karpathy/autoresearch}
\end{thebibliography}

\end{document}
